\section{Our Approach}
\label{sec:impl}

Five implementations grow one exact CART or C4.5 tree on numeric attributes.
They differ in how a node's rows stay sorted, which cuts are scored explicitly,
and where the nodes are grown.
\ImplGpu{} and \ImplHybrid{} are two settings of one CUDA builder: every node
on the GPU, or the small subtrees on the CPU.

\subsection{Common Design}
\label{sec:impl-common}

Values are stored by column as 32-bit floats, and labels become dense integers
at load time.
All five score a cut with one double-precision routine, compiled for the CPU
and the GPU.
Gains within \(10^{-12}\) are ties: the earlier cut wins, then the lower
attribute index.
A shortcut may discard a cut only when it cannot change that decision.

\ImplSerialBase{} keeps lists of row identifiers.
The other four store each attribute as 8-byte entries, the value together with
the row and the class.
A node is one index range in every column.
A stable partition of that range leaves both children sorted, so each attribute
is ordered only once, as in SLIQ and
SPRINT~\cite{mehta1996sliq,shafer1996sprint}, with the class kept in the entry.
The tree is a flat array in preorder.
Builders that create nodes concurrently write them to a shared store and number
them afterwards, so creation order does not change the tree.
The five implementations match one another.
The checks against scikit-learn and Quinlan's program are those of
Section~\ref{sec:intro}.
Pruning, when it is used, is shared CPU code after growth.
Section~\ref{sec:design} mostly reports depth-limited trees left unpruned.

\subsection{Serial Simple}
\label{sec:impl-simple}

\ImplSerialBase{} is the textbook procedure on one thread, and the reference
for the others.
For each node and attribute it copies the row identifiers, sorts them by
value, scores every admissible cut, and recurses on two new lists.
The sorts read values through those identifiers, and the lists are allocated
at every node.

\subsection{Serial Optimized}
\label{sec:impl-serial}

\ImplSerialOpt{} is the builder of Section~\ref{sec:impl-parallel} with the
thread pool off, so both the presort and the growth use one thread.
A stable radix sort orders each attribute once.
The left child is a prefix of the winning column; the other columns are stably
partitioned to match and are not sorted again.
Rows that share a value may change order without changing the class counts on
either side.
An attribute that is constant on the node is dropped for the whole subtree.

A running estimate, and a convex bound on two-class nodes, discards cuts that
cannot win under the tie rule.
The shared routine then scores the remaining cuts in order and supplies the
children's class counts.
One slot per attribute and one spare slot hold the columns for the whole run.
Large nodes partition out of place into the spare slot; smaller nodes
partition in place through a reused buffer.
Allocations of at least 2\,MiB request transparent huge pages.
An explicit stack replaces the recursion of \ImplSerialBase{}.

\subsection{Multicore CPU}
\label{sec:impl-parallel}

\ImplParallel{} uses the same columns and the same scan, and picks the
schedule by node size.
On a node of at least 32{,}768 rows, the attribute scans run together, and so
does the partition, one column per thread.
On a node of at least 4{,}096 rows, the left child is a pool task and the
current thread takes the right child.
The pool joins only at the end.
Spare slots, at most one per thread and at most one fewer than the number of
attributes, let a large node move several columns at once.
The presort uses the same pool.
The CPU builder does not pin workers or bind columns to a socket.
The pool size is the number of CPUs the process may use, so an external
\texttt{taskset} or \texttt{numactl} limit stays in force, and a page remains
on the node of the thread that first writes it.

\subsection{GPU-Based}
\label{sec:impl-gpu}

\ImplGpu{} keeps every subtree on the GPU.
The hand-off threshold is one row.
The CPU stages copies, applies the split rule once per level, and runs pruning.

All nodes of one level are processed together.
Each column range is cut into tiles, one thread block per tile, so a level is
a short sequence of launches.
The device holds the same 8-byte columns and one spare slot, and radix-sorts a
column as soon as its upload has finished.
On the A100 and the H100 of Section~\ref{sec:design}, every admissible cut is
scored in double precision with the shared routine.
Where double precision is much slower, a single-precision pass first drops
cuts that cannot win, and only the rest are scored exactly.
One best cut per attribute returns to the host.
The winning column is split as a prefix; the other columns are stably
partitioned, and that write counts the next level's tiles.

The training matrix stays pageable.
Threads copy it through a few page-locked staging buffers while earlier chunks
are in flight.
When the process may run outside the GPU's NUMA node, those buffers are
allocated on that node and the copy threads are pinned to its CPUs.
This is the only explicit NUMA placement in the source, and the copy
threads of \ImplHybrid{} are pinned in the same way.
The host training columns stay where the threads that fill them first touch
them.

\subsection{Hybrid}
\label{sec:impl-hybrid}

\ImplHybrid{} adds a size test after each GPU level.
A child with fewer than \(m\) rows is copied to the host and grown by the pool
of Section~\ref{sec:impl-parallel} while the GPU continues.
The threshold \(m\) can be set directly.
By default it is \(N/8192\), clamped to \([1024, 2048]\), for CART, and
\(N/4096\), clamped to \([1024, 4096]\), for C4.5, with \(N\) the number of
training rows.
The defaults were chosen empirically.
A level of many tiny nodes leaves the GPU with too little work per launch,
while a CPU thread can finish a small subtree without another transfer.

\begin{figure}[t]
\centering
\includegraphics[width=\linewidth]{figs/approach-hybrid.pdf}
\caption{One level of \ImplHybrid{}.
Large nodes stay on the GPU.
Children below the size threshold are copied to the host, and the CPU pool
grows them while the GPU starts the next level.
\ImplGpu{} keeps those children on the device.}
\label{fig:hybrid}
\end{figure}

Figure~\ref{fig:hybrid} shows the copy.
Small children are gathered into a device ring of fixed capacity, downloaded
through page-locked buffers, and written to a pageable host area large enough
that each row is transferred at most once.
A callback starts the CPU tasks when their columns arrive.
The GPU stalls only when the ring is full.
CPU threads touch the host area during the first levels, so the copy finds its
pages resident.
That area is not allocated on the GPU's NUMA node.
Both devices append to the store of Section~\ref{sec:impl-common}.
Each child is assigned by its row count once that count is known, and both
devices write nodes of the same tree.
